\newcommand{\irule}[3]{ \inferrule*[right=(#1)]{#2}{#3} }
\def\iaxiom#1#2{\irule{#1}{ }{#2}}

% http://tex.stackexchange.com/questions/3964/mathematical-symbol-for-disjoint-set-union
\makeatletter
\def\moverlay{\mathpalette\mov@rlay}
\def\mov@rlay#1#2{\leavevmode\vtop{%
   \baselineskip\z@skip \lineskiplimit-\maxdimen
   \ialign{\hfil$\m@th#1##$\hfil\cr#2\crcr}}}
\newcommand{\charfusion}[3][\mathord]{
    #1{\ifx#1\mathop\vphantom{#2}\fi
        \mathpalette\mov@rlay{#2\cr#3}
      }
    \ifx#1\mathop\expandafter\displaylimits\fi}
\makeatother

\newcommand{\dcup}{\charfusion[\mathbin]{\cup}{+}}
\newcommand{\bigdcup}{\charfusion[\mathbin]{\bigcup}{+}}

\def\opbin{\boxplus\xspace}
\def\opun{\boxminus\xspace}

% Types
\newcommand{\tInt}{\textsc{Int}\xspace}
\newcommand{\tFloat}{\textsc{Float}\xspace}
\newcommand{\tUnit}{\textsc{Unit}\xspace}

\def\structRec#1#2{\textsc{rec}(\&#1) #2} % TODO

\newcommand{\tFV}[1]{\textrm{FreeVars}(#1)}

\def\opxor{\mathbin{\char`\^}}

% Constant
\newcommand{\cNil}{\textsc{Nil}}

% Program & language names
\newcommand{\ctonewspeak}{\textsc{c2newspeak}}
\newcommand{\cpp}{\textsc{cpp}}
\newcommand{\newspeak}{\textsc{Newspeak}}

% Newspeak
\newcommand{\npkIf}[3]{\textsc{If}(#1) \{ #2 \} \textsc{Else} \{ #3 \}}
\newcommand{\npkDoWith}[2]{\textsc{Do}\{ #1 \}\textsc{With}~#2:}
\newcommand{\npkWhile}[1]{\textsc{While}(1)\{#1\}}
\newcommand{\npkGoto}[1]{\textsc{Goto}~#1}
\newcommand{\npkReturn}[1]{\textsc{Return}~#1}

% cML
\def\iPass{\textsc{Pass}\xspace}
\def\iIf#1#2#3{\npkIf{#1}{#2}{#3}}
\def\iWhile#1#2{\textsc{While}(#1)\{#2\}}
\def\iReturn#1{\textsc{Return}(#1)}
\def\iDecl#1#2#3{\textsc{Decl}~#1=#2~\textsc{in}\{#3\}}

\def\iThen#1#2{\textsc{If}(#1) \{ #2 \}}

\def\eNull{\textsc{Null}\xspace}
\def\eUnit{(\,)\xspace}

\def\vRet{\underline{R}}

% sets pour semantique
\def\sId{\textsc{Id}\xspace}
\def\sMem{\textsc{Mem}\xspace}
\def\sAddr{\textsc{Addr}\xspace}
\def\sInt{\textsc{Int}\xspace}
\def\sFloat{\textsc{Float}\xspace}
\def\sPhi{\Phi}

% Memory state (instructions)
% #1 : memory
% #2 : instruction
\def\msi#1#2{ \langle #2, #1 \rangle }

% Reduction (expressions)
\def\mm#1#2#3#4{ \msi{#1}{#2} → \msi{#3}{#4} }
\def\mms#1#2{\mm{m}{#1}{m}{#2}}

% Reduction (instructions)
\def\mmi#1#2#3#4{ \mm{#1}{#2}{#3}{#4} }

% Reduction (left-values)
\def\mmlv#1#2#3#4{ \mm{#1}{#2}{#3}{#4} }

% Typage sém
\def\semtyp#1#2#3{#1 ⊧ #2 : #3}
\def\semtypphi#1#2#3{#1 ⊧_Φ #2 : #3}

% Compatibilité Γ / m
\def\mcomp#1#2{#1 ⊧ #2}

% Reduction contexts
\def\ctxOp#1#2{#1~\opbin~#2}
\def\ctxEmpty{\bullet\xspace}
\def\ctxUnOp#1{\opun~#1}
\def\ctxLvField#1#2{#1.#2}
\def\ctxLvIndex#1#2{#1[#2]}
\def\ctxSet#1#2{#1 ← #2}

% Context substitution
\def\ctxSub#1#2{#1\llparenthesis#2\rrparenthesis}
\def\ctxSubL#1#2{\ctxSub{#1}{#2}_L}
\def\ctxSubI#1#2{\ctxSub{#1}{#2}_I}

% Qualifs
\def\qKernel{\textsc{Kernel}\xspace}
\def\qUser{\textsc{User}\xspace}
\def\vTainted#1{\widehat{\diamondsuit}~#1}
\def\eTaint#1{\diamondsuit~#1}

\def\ptrK#1{#1~*}
\def\ptrU#1{#1~@}

% Eval de phrase
\def\ph#1#2#3{#1 ⊩ #2 → #3}

\def\phstar#1#2#3{#1 ⊩ #2 →^* #3}

\def\phprog#1#2{⊩ #1 →^* #2}

% typage
\def\ty#1#2#3{#1 ⊢ #2 : #3}
\def\tyi#1#2{#1 ⊢ #2}
\def\typh#1#2#3{#1 ⊢ #2 → #3}
\def\typhstar#1#2#3{#1 ⊢ #2 →^* #3}

\def\sVal{\textsc{Val}}
\def\sByte{\textsc{Byte}}

\def\powerset#1{\mathcal P \left({#1}\right)}

% Functions with placeholders
\def\phx#1{\ensuremath{#1{\cdot}}}
\def\phxx#1{\ensuremath{#1{\cdot}{\cdot}}}

% Progs bien formés
\def\WF#1{\textsc{wf}(#1)}
\def\NORET#1{\textsc{NoRet}(#1)}

\newcommand{\insertcode}[1]{\input{gen/#1.pyg.tex}}

% \chapter* + add to toc
%\newcommand{\chapterstar}[1]{\chapter*{#1} \addcontentsline{toc}{chapter}{#1}}

% Register names
\def\defreg#1{\expandafter\def\csname#1\endcsname{\textsc{#1}\xspace}}

\defreg{eax}
\defreg{ebx}
\defreg{ecx}
\defreg{edx}
\defreg{esi}
\defreg{edi}
\defreg{ebp}
\defreg{esp}
\defreg{eip}
\defreg{eflags}

\def\crtrois{\textsc{CR3}}

\newcommand{\ring}[1]{\emph{ring}~#1}

\newcommand{\asminstr}[1]{\texttt{#1}}

\newcommand{\wip}[0]{%
\begin{center}
    \includegraphics[width=5cm]{gen/travaux.pdf}
\end{center}
%\textbf{(wip)}
}

\def\listcons{\texttt{(::)}\xspace}

\tikzstyle{bigbrace} = [decorate,decoration={brace,amplitude=3mm}]

% Debug tikz bounding box
\def\drawbb{\draw[thick] (current bounding box.south west) rectangle (current bounding box.north east);}

% SubFloat + verbatim
\makeatletter
\newbox\sf@box
\newenvironment{SubFloat}[2][]%
{\def\sf@one{#1}%
\def\sf@two{#2}%
\setbox\sf@box\hbox
\bgroup}%
{ \egroup
\ifx\@empty\sf@two\@empty\relax
\def\sf@two{\@empty}
\fi
\ifx\@empty\sf@one\@empty\relax
\subfloat[\sf@two]{\box\sf@box}%
\else
\subfloat[\sf@one][\sf@two]{\box\sf@box}%
\fi}
\makeatother

% comes from the memoir manual {{{

\newcommand*{\addtotextheight}[1]{%
\twocolumn
\addtolength{\textheight}{#1}%
\onecolumn}

%% for main body, bottom of text at 1in, footer below
%% top of header at 1in, first text line double spaced
%% below base of header
\newlength{\linespace}
\setlength{\linespace}{\baselineskip} % current equivalent of \onelineskip
\setlength{\headheight}{\onelineskip}
\setlength{\headsep}{\linespace}
\addtolength{\headsep}{-\topskip}
\setlength{\uppermargin}{1in}
\addtolength{\uppermargin}{\headheight}
\addtolength{\uppermargin}{\headsep}

%% and for the bottom
\setlength{\lowermargin}{1in}
\setlength{\textheight}{\paperheight}

\addtolength{\textheight}{-\uppermargin}
\addtolength{\textheight}{-\lowermargin}
%% footnote settings
\setlength{\footskip}{\onelineskip}
\setlength{\footnotesep}{\onelineskip}


%% the fiddle lengths (..ta.. for title/approval page, others for prelims)
\newlength{\toptafiddle} \setlength{\toptafiddle}{2\linespace}
\newlength{\bottafiddle} \setlength{\bottafiddle}{0pt}
\newlength{\topfiddle}
\setlength{\topfiddle}{\toptafiddle}
\newlength{\botfiddle}
\setlength{\botfiddle}{\onelineskip}

%%% general macro for Abstract, etc., headings
\newcommand*{\pretoctitle}[1]{{\clearpage\centering \vspace*{-\topfiddle}#1\par}}

%%% Start the ACKNOWLEDGEMENTS
\newcommand{\asuacknowledgements}{\pretoctitle{REMERCIEMENTS}}

%%% Start the ABSTRACT
\newcommand{\asuabstract}{ \addtotextheight{-\botfiddle} \pretoctitle{ABSTRACT}}

\addtodef{\mainmatter}{\addtotextheight{\botfiddle}{}}

%% make it easy to center any dedication
\newcommand{\asudedication}[1]{%
{\clearpage\mbox{}\vfill\centering #1 \par\vfill\clearpage}}

% }}} /memoir

% {{{ subtitle command
\newcommand\subtitle[1]{%
  \newcommand\thesubtitle{#1}}
% }}}

% Language name
\def\langname{\textsc{Safespeak}\xspace}

% Center vertically + remove header & footer
\newenvironment{headingpage}{
\thispagestyle{simple}
\vspace*{\fill}
}{
\vspace*{\fill}
}

% Hoare triple
\def\hoare#1#2#3{\ensuremath{\{#1\}~#2~\{#3\}}}

\def\eqdef{\overset{\mathrm{def}}{=\joinrel=}\xspace}

% Lens
\def\mkLens#1#2{\langle #1 | #2 \rangle}
\def\setLens#1#2{\textsc{Lens}_{#1, #2}}
\def\lensGetX#1{\ensuremath{\mathrm{get}_{#1}}}
\def\lensPutX#1{\ensuremath{\mathrm{put}_{#1}}}
\def\lensGet#1#2{\lensGetX{#1}(#2)}
\def\lensPut#1#2#3{\lensPutX{#1}(#2, #3)}
\def\lensNodeBig#1#2{
\vcenter{\hbox{
\begin{tikzpicture}[scale=0.7]
\draw[fill=#1] (-0.5,-0.5) rectangle ++ (1,1);
\draw[fill=#2] (0,0) circle (0.3);
\end{tikzpicture}
}}
}
\def\lensInner#1{
\vcenter{\hbox{
\begin{tikzpicture}[scale=0.7]
\draw[fill=#1] (0, 0) circle (0.3);
\end{tikzpicture}
}}
}
\def\lensNodeBigIdx#1#2#3#4#5{
\vcenter{\hbox{
\begin{tikzpicture}[scale=0.7]
  \draw[fill=#1] (-1.5,-1.5) rectangle ++ (3,3);
  \node[fill=#2, draw, star, minimum size=6mm] at (-0.75, 0.75) {};
  \node[fill=#3, draw, star, minimum size=6mm] at ( 0.75, 0.75) {};
  \node[fill=#4, draw, circle, minimum size=6mm] at (-0.75,-0.75) {};
  \node[fill=#5, draw, regular polygon, regular polygon sides=5, minimum size=6mm] at ( 0.75,-0.75) {};
  \node at (-2,  0.75) {a};
  \node at ( 2,  0.75) {b};
  \node at (-2, -0.75) {c};
  \node at ( 2, -0.75) {d};
\end{tikzpicture}
}}
}
\def\lensInnerStar#1{
\vcenter{\hbox{
\begin{tikzpicture}[scale=0.7]
\node[draw, fill=#1, star] {};
\end{tikzpicture}
}}
}

% Erreurs

\def\serr#1{
  \switch{#1}
    {taint}{Ω_{taint}}
    {array}{Ω_{array}}
    {ptr}{Ω_{ptr}}
    {div}{Ω_{div}}
    {field}{Ω_{field}}
    {var}{Ω_{var}}
    {END}
    % important : étendre la description dans le chapitre 4 et la figure de
    % liste
}

%\def\uGet#1#2{#1 ⇐_U #2}
%\def\uPut#1#2{#1 ⇒_U #2}
\def\cfuF{\textrm{copy\_from\_user}\xspace}
\def\ctuF{\textrm{copy\_to\_user}\xspace}
\def\uGet#1#2{\cfuF(#1, #2)}
\def\uPut#1#2{\ctuF(#1, #2)}

% Litteral notation

\def\eStruct#1{\{#1\}}
\def\eArray#1{[#1]}
\def\tStruct#1{\eStruct{#1}}
\def\eFun#1#2{\mathrm{fun} (#1) \{ #2 \}}

\def\gramorx{\hspace{3mm}\gramor\xspace}
\def\stFun#1{\textsc{Fun}_n}
\def\tComp#1#2{\mathrm{Comp}(#1, #2)}

% Box on top on rules fig
\def\ruleheader#1{
  \begin{minipage}{\linewidth}
    \raggedleft
    \fbox{#1}
  \end{minipage}
}

\def\termFst#1{\ensuremath{\textrm{fst}~#1}}
\def\termSnd#1{\ensuremath{\textrm{snd}~#1}}

\def\ptrtype{\texttt{ptrtype}\xspace}

\def\addrInt{\textrm{addr}^\natural \xspace}
\def\interpPtr#1{\textrm{interp}_{(#1)}}

\tikzset{>=latex}

\def\partccl#1{
  \clearpage~\thispagestyle{cleared}\cleartorecto
  \chapter*{Conclusion de la partie~\ref{#1}}
  \addcontentsline{toc}{chapter}{Conclusion de la partie~\ref{#1}}
}

\def\hexa#1{0x#1}
\def\floatop#1{#1.}

\def\cLookup#1#2{\mathrm{Lookup} (#1, #2)}
\def\cExtend#1#2#3{\mathrm{Extend}(#1, #2 ↦ #3)}
\def\cCleanup#1{\mathrm{Cleanup}(#1)}

\newcommand{\motif}[1]{\cleaders\hbox{#1}\hfil}
\newcommand{\chain}[2]{\leavevmode\hbox to #2{\motif{#1}}}
\newcommand{\diamonds}{$\mkern-.6mu \diamond \mkern-.6mu$}

\def\jolibreak{
\fancybreak{\chain{\diamonds}{0.25\textwidth}}
}

\def\cDom#1{\mathrm{dom}(#1)}

\def\cVarname#1{\mathrm{Varname}(#1)}

\def\cVisible#1{\mathrm{Visible}(#1)}

\def\cRepr#1{\mathrm{Repr}(#1)}

\def\cPush#1#2{\mathrm{Push}(#1, #2)}

\def\cPop#1{\mathrm{Pop}(#1)}

\def\cH#1{\mathcal{H}(#1)}
\def\cHphi#1{\mathcal{H}_Φ(#1)}
\def\cL#1#2{\mathrm{CleanL_{#1}(#2)}}
\def\cV#1#2{\mathrm{CleanV_{#1}(#2)}}
